\subsection{LIE GROUP CASE}

PROPOSITION 6. Assume that k is R, C or a non-discrete complete ultrametric field of characteristic 0. Let G be a finite dimensional Lie group over k, e its identity element, g its Lie algebra, h a Cartan subalgebra of g, \(h_{r}\) the set of regular elements of g belonging to h.

\begin{enumerate}
\def\labelenumi{(\roman{enumi})}
\item
  Let \(\mathfrak{s}\) be a vector space complement of \(\mathfrak{h}\) in \(\mathfrak{g}\) , \(\mathfrak{s}_0\) a neighbourhood of 0 in \(\mathfrak{s}\) on which an exponential map is defined, and \(h_0 \in \mathfrak{h}_r\) . The map \((s, h) \mapsto \mathrm{F}(s, h) = (\exp \operatorname{ad} s).h\) from \(\mathfrak{s}_0 \times \mathfrak{h}\) to \(\mathfrak{g}\) is étale at \((0, h_0)\) .
\item
  The map \((g,h)\mapsto \mathrm{F}'(g,h) = (\operatorname {Ad}g).h\) from \(\mathbf{G}\times \mathfrak{h}_r\) to \(\mathfrak{g}\) is a submersion. In particular, its image \(\Omega\) is open. For all \(x\in \Omega\) , \(\mathfrak{g}^0 (x)\) is a Cartan subalgebra of \(\mathfrak{g}\) conjugate to \(\mathfrak{h}\) under \(\operatorname {Ad}(\mathrm{G})\)
\item
  Let \(h_0 \in \mathfrak{h}_r\) . For any neighbourhood \(\mathrm{U}\) of \(e\) in \(\mathrm{G}\) , the set \(\bigcup_{a \in \mathrm{U}} (\mathrm{Ad}a)(\mathfrak{h}_r)\) is a neighbourhood of \(h_0\) in \(\mathfrak{g}\) .
\end{enumerate}

Let \(h_0\) and \(\mathfrak{s}\) be as in (i). Let \(\mathrm{T}\) be the tangent linear map of \(\mathrm{F}\) at \((0, h_0)\) . Then \(\mathrm{F}(0, h) = h\) for all \(h \in \mathfrak{h}\) , so \(\mathrm{T}(0, h) = h\) for all \(h \in \mathfrak{h}\) . On the other hand, for \(\mathfrak{s}_0\) sufficiently small, the tangent linear map at 0 of the map \(s \mapsto \exp s\) from \(\mathfrak{s}_0\) to \(\operatorname{End}(\mathfrak{g})\) is the map \(s \mapsto \operatorname{ad}s\) from \(\mathfrak{s}\) to \(\operatorname{End}(\mathfrak{g})\) . Thus \(\mathrm{T}(s, 0) = [s, h_0]\) for all \(s \in \mathfrak{s}\) . Now the map from \(\mathfrak{g}/\mathfrak{h}\) to \(\mathfrak{g}/\mathfrak{h}\) induced by \(\operatorname{ad}h_0\) by passage to the quotient is bijective. It follows that \(\mathrm{T}\) is bijective, hence (i). Since \(\exp s = \operatorname{Ad}\exp s\) for all \(s \in \mathfrak{s}\) sufficiently close to 0, (iii) and the first assertion of (ii) follow. Every \(x \in \Omega\) is of the form \((\operatorname{Ad}a)(h)\) with \(a \in G\) and \(h \in \mathfrak{h}_r\) , so \(\mathfrak{g}^0(x) = (\operatorname{Ad}a)(\mathfrak{g}^0(h)) = (\operatorname{Ad}a)(\mathfrak{h})\) is a subalgebra of \(\mathfrak{g}\) conjugate to \(\mathfrak{h}\) under \(\operatorname{Ad}(G)\) .

